\chapter{Conclusions and Future Work} \label{chap:conclusion}
\section{Achievements}
\section{Future Work}
\section{Final Remarks}
This thesis began with the observation of a critical gap in the MIR research community: the absence of a unified framework for generating time-aligned, multi-modal visualizations suited to contemporary research workflows. Through an iterative development process and careful architectural design, we have delivered LayerIt, a Python-based, object-oriented tool that addresses this gap while explicitly supporting future extension and integration.

LayerIt’s value lies not merely in its immediate applicability to beat tracking and annotation tasks, but in its demonstration that visualization frameworks can be architected for extensibility, composability, and community adoption. By grounding the system in established conventions (OO design patterns, common MIR libraries, standardized formats), we have created a foundation upon which the community can build up on. 

As MIR research continues to evolve, driven by advances in machine learning, increasing dataset complexity, and growing interdisciplinary collaboration. The need for sophisticated, flexible visualization tools will only intensify. LayerIt is positioned to serve as a foundation for this future work, enabling researchers to focus on their core innovations rather than reimplementing visualization infrastructure from scratch.
